\subsection{指数映射}

定理 4. 设 G 是李群芽，L 是其李代数，V 是 L 中 0 的一个开邻域，\(\phi\) 是从 V 到 G 的解析映射，满足 \(\phi(0) = 0\) 及 \(\mathrm{T}_0(\phi) = \mathrm{Id}_L\)。以下条件等价：

\begin{enumerate}
\def\labelenumi{(\roman{enumi})}
\item
  对所有 \(b \in \mathbf{L}\)，有 \(\phi((\lambda + \lambda')b) = \phi(\lambda b)\phi(\lambda'b)\)，其中 \(|\lambda|\) 和 \(|\lambda'|\) 充分小。
\item
  对所有 \(b \in \mathbf{L}\) 及每个整数 \(n > 0\)，\(\phi_{*}(b^{n})\) 是 \(n\) 次齐次元素，位于

\(\mathrm{U}(\mathrm{G})\) 中（将 \(\mathrm{T}_{0}^{(\infty)}(\mathrm{L})\) 与 \(\mathrm{TS}(\mathrm{L})\) 识别，并把 \(b^{n}\) 视为在 \(\mathrm{TS}(\mathrm{L})\) 中计算）。

\item
  TS(L) 到 U(G) 的映射 \(\phi_{*}\) 与 TS(L) 和 U(G) 的分次相容。
\item
  TS(L) 到 U(G) 的映射 \(\phi_*\) 是 TS(L) 到 L 的包络代数的典范映射。
\item
  存在一个定义 \(\mathbf{L}\) 的拓扑的范数，并且使 \(\mathbf{L}\) 满足
\end{enumerate}

\[
\| [ x, y ] \| \leqslant \| x \| \| y \|
\]

对所有 \(x, y\)，其中元素属于 \(\mathbf{L}\)，并且存在一个开子群芽 \(\mathbf{W} \subset \mathbf{V}\)，它是由 \(\mathbf{L}\) 定义的李群芽（第2号）的开子群芽，使得 \(\phi | \mathbf{W}\) 是从 \(\mathbf{W}\) 到 \(\mathbf{G}\) 的一个开李子群芽的同构。

(v) \(\Rightarrow\) (i)：显然，因为在 W 中有 \((\lambda b).(\lambda^{\prime}b) = (\lambda +\lambda^{\prime})b\)，其中 \(|\lambda |\) 和 \(|\lambda '|\) 充分小。
(i) \(\Rightarrow\) (ii)：假设条件 (i) 成立。设 \(b \in L\)。设 \(\psi\) 是 \(\phi\) 限制在 \(V \cap Kb\) 上的映射。根据假设，存在 0 的一个对称邻域 \(T\)，它位于加法李群 \(Kb\) 中，使得 \(\psi|T\) 是从李群芽 \(T\) 到 \(G\) 的态射。因此

\[
\phi_ {*} (b ^ {n}) = (\psi | \mathrm{T}) _ {*} (b ^ {n}) = ((\psi | \mathrm{T}) _ {*} (b)) ^ {n} = (\phi_ {*} (b)) ^ {n},
\]

所以 \(\phi_{*}(b^{n})\) 是 \(n\) 次齐次元素，属于 U(G)。

(ii) \(\Rightarrow\) (iii)：这是因为 \(\mathsf{TS}^n (\mathbf{L})\) 是 \(\mathsf{TS}(\mathbf{L})\) 的向量子空间，由 \(n\) 次幂（取自 \(\mathbf{L}\) 中元素）生成（《代数》第 IV 章 § 5 命题5）。
(iii) \(\Rightarrow\) (iv)：TS(L) 到 L 的包络代数的典范映射，是将 1 映到 1 且延拓 \(\mathrm{Id}_{\mathrm{L}}\) 的分次余代数的唯一态射（第 II 章 § 1，第5号，注记3）。现在，根据假设，\(\phi_{*}\) 是余代数态射且 \(\phi_{*}|_{\mathrm{L}} = \mathrm{Id}_{\mathrm{L}}\)。若条件 (iii) 成立，容易看出条件 (iv) 也成立。
(iv) \(\Rightarrow\) (v)：假设条件 (iv) 成立。在 L 上取一个定义 L 的拓扑的范数，使 \(\| [x,y]\| \leqslant \| x\| \| y\|\) 对所有 \(x,y\) 成立。设 H 是由赋范李代数 L 定义的李群芽。根据定理1，存在 H 的一个开子群芽 \(S\subset V\) 以及从 S 到 G 的一个开子群芽的同构 \(\phi^{\prime}\)。由于我们已经知道 (v) \(\Rightarrow\) (iv)，TS(L) 到 U(G) 的映射 \(\phi_{*}\) 是 TS(L) 到 L 的包络代数的典范映射。因此有 \(\phi_{*}(t) = \phi_{*}^{\prime}(t)\)，对所有 \(t\in T_0^{(\infty)}(L)\)。由于 \(\phi\) 和 \(\phi^{\prime}\) 都解析，\(\phi\) 与 \(\phi^{\prime}\) 在 0 的某个邻域上重合。

定义 1. 设 G 是李群芽，L 是其李代数。G 的指数映射是定义在 L 中 0 的一个开邻域上、取值于 G 且满足定理4条件的任意解析映射 \(\phi\)。

定理4立即推出：对每个李群芽 G，都存在 G 的指数映射，并且 G 的两个指数映射在 0 的某个邻域上重合。

例。(1) 设 G 是完备可赋范空间 E 的加法群。L(G) 到 E 的典范同构满足定理4的条件 (i)，因而是 G 的指数映射。

\begin{enumerate}
\def\labelenumi{(\arabic{enumi})}
\setcounter{enumi}{1}
\tightlist
\item
  设 A 是完备赋范含幺结合代数。设 \(\mathbf{A}^*\) 是由 A 的可逆元素组成的李群。我们将 \(\mathrm{L}(\mathrm{A}^*)\) 与 A 识别（§3，第9号，命题33的推论）。若 \(\mathbf{K} = \mathbf{R}\) 或 \(\mathbf{C}\)，我们知道第 II 章 §7 第3号中定义的从 A 到 \(\mathbf{A}^*\) 的映射 exp 满足定理4的条件 (i)，因而是指数映射。现在设 K 是超度量的。设 \(p\) 是 K 的剩余域的特征。若 \(p \neq 0\)，令 \(\lambda = |p|^{1/(p-1)}\)；若 \(p = 0\)，令 \(\lambda = 1\)。设 U 是 \(x \in \mathbf{A}\) 且满足 \(\|x\| < \lambda\) 的集合。我们知道（第 II 章 §8，第4号），从 U 到 \(\mathbf{A}^*\) 的映射 exp 满足定理4的条件 (i)，因而是指数映射。注意 U 是 A 的加法子群。
\end{enumerate}

这个例子说明了定义1所采用术语的由来。

设 \(G\) 是李群芽，\(\phi\) 是 \(G\) 的指数映射。则 \(\phi\) 在 0 处是 étale 的，因而存在一个开邻域 \(U\)，它位于 \(L(G)\) 中的 0 附近，使得 \(\phi(U)\) 在 \(G\) 中开，且 \(\phi|U\) 是从解析流形 \(U\) 到解析流形 \(\phi(U)\) 的同构。

G 上的典范图（第一类）是解析流形 G 上的一个图 \(\psi\)，其逆映射是指数映射。若进一步 G 是有限维的，并选定 L(G) 的一组基，则由 \(\psi\) 和该基在 \(\psi\) 的定义域上定义的坐标系称为典范坐标系（第一类）。

命题 3. 设 \(\mathbf{G}\) 是李群芽，\(\mathbf{L}\) 是其李代数，\(\phi\) 是 \(\mathbf{G}\) 的指数映射。设 \(\mathbf{L}_1, \ldots, \mathbf{L}_n\) 是 \(\mathbf{L}\) 的向量子空间，使得 \(\mathbf{L}\) 是 \(\mathbf{L}_1, \ldots, \mathbf{L}_n\) 的拓扑直和。映射

\[
(b _ {1}, b _ {2}, \dots , b _ {n}) \mapsto \theta (b _ {1}, b _ {2}, \dots , b _ {n}) = \phi (b _ {1}) \phi (b _ {2}) \dots \phi (b _ {n}),
\]

定义在 \(L_{1} \times L_{2} \times \cdots \times L_{n}\) 的一个开子集上，是解析的。在 \((0, 0, \ldots, 0)\) 处，\(\theta\) 的切映射是从 \(L_{1} \times \cdots \times L_{n}\) 到 L 的典范映射。

设 \(k_{i}\) 是从 \(L_{i}\) 到 \(L_{1} \times L_{2} \times \cdots \times L_{n}\) 的典范嵌入。则对所有 \(b \in L_{i}\)，\((\mathrm{T}_{(0,\ldots,0)}\theta)(\mathrm{T}_{0}\phi_{i})(b) = (\mathrm{T}_{0}\phi)(b) = b\)，从而 \((\mathrm{T}_{(0,\ldots,0)}\theta)|L_{i}\) 是从 \(L_{i}\) 到 L 的典范嵌入。

特别地，\(\theta\) 在 \((0, 0, \ldots, 0)\) 处是 étale 的。它限制在 \((0,0,\ldots,0)\) 的充分小开邻域 U 上时，在 G 中的像 \(\theta(\mathrm{U})\) 是开的，且这个限制是一个同构。逆映射 \(\eta\) 把 \(\theta(\mathrm{U})\) 映到 U，称为 G 的第二类典范图，它与给定的 L 的直和分解相关。若进一步 G 是有限维的，且每个 \(L_{i}\) 都由非零向量 \(e_{i}\) 生成，则在 \(\theta(\mathrm{U})\) 上由 \(\eta\) 和 \(e_{i}\) 定义的坐标系称为第二类典范坐标系。

命题 4. 设 G 是李群芽，\(\phi\) 是 G 的单射指数映射。对 L(G) 中的所有 \(x, y\)，


\[
x + y = \lim _ {\lambda \in \mathbb {K} ^ {*}, \lambda \rightarrow 0} \lambda^ {- 1} \phi^ {- 1} (\phi (\lambda x) \phi (\lambda y))\tag{1}
\]

\[
[ x, y ] = \lim _ {\lambda \in \mathbb {K} ^ {*}, \lambda \rightarrow 0} \lambda^ {- 2} \phi^ {- 1} (\phi (\lambda x) \phi (\lambda y) \phi (- \lambda x) \phi (- \lambda y))\tag{2}
\]

（注意，\(\phi^{-1}(\phi(\lambda x)\phi(\lambda y))\) 和 \(\phi^{-1}(\phi(\lambda x)\phi(\lambda y)\phi(-\lambda x)\phi(-\lambda y))\) 在 \(|\lambda|\) 充分小时有定义。）

给 \(\mathbf{L} = \mathbf{L}(\mathbf{G})\) 赋予一个定义 \(\mathbf{L}\) 拓扑的范数，使 \(\| [x,y] \| \leqslant \| x \| \| y \|\) 对所有 \(x, y\) 属于 \(\mathbf{L}\) 成立。利用定理2和定理4，可以假设 \(\mathbf{G}\) 是由 \(\mathbf{L}\) 定义的李群芽，且 \(\phi = \mathrm{Id}_{\mathbf{G}}\)。记 \((x,y) \mapsto x.y\) 为群 \(\mathbf{G}\) 中的乘积。于是待证公式可写为


\[
x + y = \lim _ {\lambda \in \mathbf {K} ^ {*}, \lambda \rightarrow 0} \lambda^ {- 1} ((\lambda x). (\lambda y))\tag{3}
\]

\[
[ x, y ] = \lim _ {\lambda \in \mathbb {K} ^ {*}, \lambda \rightarrow 0} \lambda^ {- 2} ((\lambda x). (\lambda y). (- \lambda x). (- \lambda y)).\tag{4}
\]

存在 K 中 0 的一个开邻域 V，使函数

\[
\lambda \mapsto f (\lambda) = (\lambda x). (\lambda y)
\]

在 V 上有定义且解析。根据第 II 章 § 6 第4号注记2，\(f\) 在原点处的积分级数展开为

\[
\lambda (x + y) + \frac {1}{2} \lambda^ {2} [ x, y ] + \dots
\]

这证明了 (3)。另一方面，对于 G 中的 \(u, v\)，当 \(\| u \|\)、\(\| v \|\) 充分小时，\(u.v\) 是 \((u, v)\) 的解析函数，而这个函数在原点处的积分级数展开中的 1 次和 2 次项为 \(u + v + \frac{1}{4}[u, v]\)。根据《可微与解析流形》R，3.2.7 和 4.2.3，函数 \(f(\lambda)f(-\lambda)\) 在原点处的积分级数展开中的 1 次和 2 次项，就是

\[
f (\lambda) + f (- \lambda) + \frac {1}{2} [ f (\lambda), f (- \lambda) ]
\]

或者也可写为

\[
\begin{array}{r l} \lambda (x + y) + \frac {1}{2} \lambda^ {2} [ x, y ] - \lambda (x + y) + \frac {1}{2} \lambda^ {2} [ x, y ] & \\ + \frac {1}{2} [ \lambda (x + y), - \lambda (x + y) ] = \lambda^ {2} [ x, y ] \end{array}
\]

这证明了 (4)。

命题 5. 设 G 是李群，k 是 K 的非离散闭子域，G' 是将 G 看作定义在 k 上的李群所得的群，且 \(\phi\)（分别为 \(\phi'\)）是 G（分别为 G'）的指数映射。则 \(\phi\) 与 \(\phi'\) 在 0 的某个邻域上重合。

\(\phi\) 相对于 \(G'\) 满足定理4的假设 (i)，因而是 \(G'\) 的指数映射。

命题 6. 设 G 是李群芽，L 是其李代数，\(\phi: V \to G\) 是 G 的指数映射。对任意 \(x \in V\)，将 \(\mathrm{T}_x(L)\) 与 L 识别，从而 \(\varpi(x)\) 是 \(\phi\) 在 \(x\) 处的右微分，是从 L 到 L 的线性映射。当 \(x\) 充分接近 0 时，

\[
\varpi (x) = \sum_ {n \geqslant 0} \frac {1}{(n + 1) !} (\mathrm{ad} x) ^ {n}.
\]

给 \(\mathbf{L}\) 赋予一个与其拓扑相容且满足 \(\| [x,y]\| \leqslant \| x\| \| y\|\)（对所有 \(x,y\in L\)）的范数。只需考虑 \(G\) 是由 \(L\) 定义的李群芽且 \(\phi = \mathrm{Id}_{\mathrm{G}}\) 的情形。按定义，此时 \(\varpi (x)\) 是在 \(x\) 处的映射 \(y\mapsto y.x^{-1}\) 的切映射，该映射从 \(G\) 到 \(G\)。若 \(H(X,Y)\) 表示 Hausdorff 级数，则 \(\varpi (x)\) 在 \(\| x\|\) 充分小时，是在 0 处的映射 \(y\mapsto H(x + y, - x)\) 的切映射，该映射从 \(G\) 到 \(G\)。在 \(H(X + Y, - X)\) 中，关于 \(Y\) 的一次项之和为

\[
\sum_ {m \geqslant 0} \frac {1}{(m + 1) !} (\mathrm{adX}) ^ {m} \mathrm{Y}
\]

（第 II 章 § 6，第5号，命题5）。命题于是由《可微与解析流形》R，3.2.4 和 4.2.3 推出。

设 G 是李群芽且 \(t \in K\)。G 的 t 次幂映射是这样一个映射：它在 e 的某个开邻域上定义且解析，取值于 G，并在 e 的某个邻域上与映射

\[
g \mapsto \phi (t \phi^ {- 1} (g))
\]

重合，其中 \(\phi\) 是 G 的单射指数映射。

命题 7. (i) 若 \(t \in \mathbf{Z}\)，则 \(t\) 次幂映射在 \(e\) 的某个邻域上与映射 \(g \mapsto g^t\) 重合。

\begin{enumerate}
\def\labelenumi{(\roman{enumi})}
\setcounter{enumi}{1}
\item
  在 \(e\) 处，\(t\) 次幂映射的切映射是比值为 \(t\) 的伸缩。
\item
  若 \(h\) 是 \(t\) 次幂映射，\(h'\) 是 \(t'\) 次幂映射，定义在 \(G\) 上，则 \(h \circ h'\) 是 \((tt')\) 次幂映射，而 \(g \mapsto h(g)h'(g)\) 是 \((t + t')\) 次幂映射。
\item
  若 \(h\) 是 \(t\) 次幂映射且 \(u \in \mathbf{U}^n(\mathbf{G})\)，则 \(h_*(u) = t^n u\)。
\end{enumerate}

只需在 G 是由完备赋范李代数定义的李群芽，且所考虑的 t 次幂映射使用指数映射 \(\phi = \mathrm{Id}_{\mathrm{G}}\) 构造时证明命题。但在这种情形下一切都是显然的。

